% Old stub, kept for reference:
%
% \section{Appendix}
% \subsection{Usage instructions}
%
% parameters, CLI flags, test case flags, tests, propagator, input, output, plotting (plot\_slice documentation), diagnostic dashboard, profile dashboard, cluster scripts, etc.
%
% careful with divb diagnostic in B=0 areas (e.g. Loop)

\section{Usage instructions}

This appendix is a quick guide for running the MHD part of SPH-EXA. It assumes familiarity with the hydro side of the code: building, launching, glass blocks, HDF5 dumps and restarts all work exactly as before. Only what is specific to MHD is described here. The code is available in the \href{https://github.com/MuellerNico/sphexa/}{\textbf{GitHub fork}}.

The MHD extension adds two new propagators (\texttt{magneto-ve} and \texttt{magneto-turb-ve}), a set of MHD test cases, additional command-line flags, and a group of particle fields that are written
to the HDF5 dumps under a \texttt{magneto::} prefix. Nothing else in the workflow changes.

\subsection{Running a simulation}

The same executable is used as for pure hydrodynamics
(\texttt{sphexa} on CPU, \texttt{sphexa-cuda} on GPU). A minimal run:

\begin{verbatim}
sphexa --init orszag-tang --prop magneto-ve --glass 30c.h5 \
       -n 240 -s 1.0 -w 0.1
\end{verbatim}

The propagator must be an MHD propagator. Use \texttt{-{}-prop magneto-ve} for all test cases, or \texttt{-{}-prop magneto-turb-ve} for driven turbulence. Using a different propagator with an MHD test will crash the program. Same goes when using an MHD propagator with a hydro test which does not allocate a magnetic field. 

The code works in units with $\mu_0 = 1$, so the magnetic pressure is $B^2/2$ and the
Alfv\'en speed is $v_A = |\mathbf{B}|/\sqrt{\rho}$.

\subsection{Test cases}

The native MHD test cases are selected with \texttt{-{}-init}. Most of them require a glass
block (\texttt{--glass}). All boxes are fully periodic. The quasi-2D tests use a thin slab in $z$ whose thickness is set automatically from the resolution.

\begin{table}[htbp]
\centering
\small
\setlength{\tabcolsep}{4pt}
\begin{tabularx}{\linewidth}{@{}lX@{}}
\toprule
\texttt{-{}-init} & description \\
\midrule
\texttt{alfven-wave}              & Circularly polarized Alfv\'en wave travelling obliquely through the box. Standard convergence test. \\
\texttt{brio-wu}                  & MHD shock tube \\
\texttt{mhd-loop}                 & Advection of a weak magnetic field loop \\
\texttt{orszag-tang}              & Orszag--Tang vortex \\
\texttt{mhd-rotor}                & Rotating dense disk in a uniform field \\
\texttt{sedov-magneto}            & Sedov blast wave with a uniform background field along $(\hat{x}+\hat{z})/\sqrt{2}$. \\
\texttt{kelvin-helmholtz-magneto} & Kelvin--Helmholtz instability with a uniform field along $x$. \\
\texttt{turbulence-magneto}       & Driven turbulence with a uniform mean field along $z$. Requires \texttt{-{}-prop magneto-turb-ve}. \\
\bottomrule
\end{tabularx}
\caption{MHD test cases.}
\end{table}

The resolution flag \texttt{-n} keeps its usual meaning of the effective number of particles per box length for all glass-based cases. For cubic boxes (Sedov, turbulence) this results in $n^3$ particles. The $z$-invariant slabs only use a single glass block in thickness. Therefore, it is worth using a smaller block (e.g. \texttt{30c.h5} or smaller) with these to not waste too many particles in $z$ direction. KH and Brio-Wu are exceptions regarding use of $-n$, due to density contrasts. Check the init files for specific documentation.

Just like with the hydro side, test-case parameters (physical constants) can be overridden by passing an HDF5 settings file after a colon (e.g. \texttt{-{}-init mhd-loop:settings.h5}).

\subsection{Command-line options}
The MHD extension adds the \texttt{-{}-resistivity} flag. It accepts:

\begin{center}
\begin{tabular}{@{}l p{0.72\linewidth}@{}}
\texttt{switch} & Tricco \& Price (2013) switch\\
\texttt{SLR} (default) & Slope-limited reconstruction of $\mathbf{B_{ab}}$\\
\texttt{SLRB} / \texttt{SLRB2} & As \texttt{SLR}, with an additional Balsara-like modulation (power 1 and 2 respectively)\\
number & Constant $\alpha_B$ (e.g. \texttt{-{}-resistivity 0.5})\\
\end{tabular}
\end{center}

Use \texttt{sphexa -{}-help} to list all available CLI flags.


\subsection{Output and restart}

The magnetic fields are written into the same HDF5 dumps as the hydrodynamic ones, under
the prefix \texttt{magneto::} (e.g.\ \texttt{magneto::Bx}). On the command line the fields
are named \emph{without} the prefix:

\begin{verbatim}
-f x,y,z,h,m,rho,vx,vy,vz,Bx,By,Bz,divB
\end{verbatim}

If \texttt{-f} is omitted, all conserved fields are written and the dump is restartable.
The conserved magnetic fields are:

\begin{verbatim}
Bx, By, Bz, dBx, dBy, dBz, dBx_m1, dBy_my, dBz_m1, psi_ch, d_psi_ch, d_psi_ch_m1
\end{verbatim}

If you pass a custom \texttt{-f} list and want to restart from the dump, all of these must be
included. Other, non-conserved (diagnostic) fields are:

\begin{verbatim}
divB, curlB_x, curlB_y, curlB_z, gradB_norm, alpha_B, du_diss
\end{verbatim}

\subsection{Diagnostics}

As in hydrodynamics, one line per iteration is appended to \texttt{constants.txt}. MHD runs add four columns. In order:

\begin{verbatim}
iteration  ttot  minDt  etot  ecin  eint  egrav  linmom  angmom
           eMag  meanDivBError  maxDivBError  resHeating  [test-specific]
\end{verbatim}

\texttt{eMag} is the total magnetic energy and is included in \texttt{etot}, so energy
conservation can be monitored directly. The two divergence columns report the
divergence error $\varepsilon_\mathrm{divB} = h\,|\nabla\cdot\mathbf{B}|/|\mathbf{B}|$, averaged over all particles and as a global maximum. \texttt{resHeating} is the instantaneous total heating rate produced by artificial resistivity. The trailing column is present only for the test cases that define one (RMS Mach number for turbulence, mode amplitude for Kelvin--Helmholtz).

\subsection{Plotting}
\todo{revise plotting subsection}
Two scripts in \texttt{scripts/} cover most needs. Both are plain Python
(\texttt{numpy}, \texttt{h5py}, \texttt{matplotlib}) and are usually called at the end of
a job script.

\subsubsection*{\texttt{plot\_constants.py} --- diagnostics dashboard}

\begin{verbatim}
python scripts/plot_constants.py out/123456/constants.txt
\end{verbatim}

Produces \texttt{constants\_diagnostics.pdf} next to the input file: energy components,
relative total-energy drift, linear and angular momentum, magnetic energy normalized to
its initial value, the mean and maximum $\nabla\cdot\mathbf{B}$ error, the time step, and
the resistive heating rate together with its cumulative integral. The magnetic-energy
panel also shows the dissipation-only expectation
$1 - \int \dot{E}_\mathrm{res}\,\mathrm{d}t / E_\mathrm{mag,0}$; the gap between the two
curves is magnetic energy created by the ideal induction term. Use
\texttt{--kind kh} to label the trailing column of a Kelvin--Helmholtz run correctly
(default \texttt{turb}).

\subsubsection*{\texttt{plot\_slice.py} --- 2D slices}

Renders an SPH-interpolated slice of any field through the box:

\begin{verbatim}
python scripts/plot_slice.py dump.h5 -i                 # metadata + field list
python scripts/plot_slice.py dump.h5 --field rho        # final step
python scripts/plot_slice.py dump.h5 5 --field Bmag -l  # step 5, log scale
python scripts/plot_slice.py dump.h5 --all --field Bmag -j 8
\end{verbatim}

A file may be followed by a step number (negative counts from the end, absent means the
last step). Several files are tiled into a single figure sharing one color scale, which is
the standard way to compare resolutions or dissipation schemes:

\begin{verbatim}
python scripts/plot_slice.py a.h5 b.h5 c.h5 --labels 128 256 512 --clean
\end{verbatim}

The most useful options are \texttt{--field}, \texttt{--axis}/\texttt{--pos} (slice plane;
by default the box midplane normal to $z$), \texttt{-r} (grid resolution),
\texttt{--cmap}, \texttt{-l} (log scale), \texttt{--vmin}/\texttt{--vmax},
\texttt{--fieldlines} (streamplot overlay, using $\mathbf{B}$ by default),
\texttt{--smooth} (widen the rendering kernel for a smoother picture),
\texttt{--scatter} (skip the interpolation, fast preview), \texttt{-j} (parallel workers,
pass \texttt{"\$SLURM\_CPUS\_PER\_TASK"} in a job script) and \texttt{--clean} (PDF output
without title for use as a figure). \texttt{-h} lists the rest.

Besides the raw datasets, the script understands derived fields computed on the fly:

\begin{table}[h]
\centering
\begin{tabular}{@{}ll@{}}
\toprule
name & quantity \\
\midrule
\texttt{Bmag}, \texttt{vmag}       & $|\mathbf{B}|$, $|\mathbf{v}|$ \\
\texttt{Emag}, \texttt{Pmag}       & magnetic energy density, magnetic pressure \\
\texttt{KE}                        & kinetic energy density \\
\texttt{curlBmag}                  & $|\nabla \times \mathbf{B}|$ \\
\texttt{divBerr}, \texttt{log\_divBerr} & $h\,|\nabla\cdot\mathbf{B}|/|\mathbf{B}|$ (and its $\log_{10}$) \\
\texttt{Bphi}, \texttt{Bperp}, \texttt{Bperp\_rel} & loop-frame decomposition for the \texttt{mhd-loop} test \\
\bottomrule
\end{tabular}
\caption{Derived fields available to \texttt{plot\_slice.py}. \texttt{-i} lists those
resolvable from a given dump.}
\end{table}

Both scripts only need the dump and are independent of the simulation build, so they can
be run on a login node or locally.

\subsection{Caveats}
\todo[inline]{move parts to appendix B}
\begin{itemize}
    \item \textbf{The $\nabla\cdot\mathbf{B}$ diagnostic is unreliable for tests with $\mathbf{B} \approx 0$ regions.} The divergence error is normalized by $|\mathbf{B}|$. Particles with exactly $\mathbf{B} = 0$ are excluded, but near-zero ones are not. So in regions with $\mathbf{B} \approx 0$ it is dominated by noise and can be arbitrarily large and corrupt both the maximum and mean divergence error. This affects in particular \texttt{mhd-loop}. You can instead read the unnormalized \texttt{magneto::divB} from the dump. 
    \item \textbf{Divergence-cleaning parameters are fixed} at $f_\mathrm{clean} = \sigma_c = 1$ and are not exposed on the command line. A more aggressive cleaning speed would require an additional Courant limit on the time step.
    \item \textbf{No block time stepping.} There is no MHD analogue of the \texttt{ve-bdt} propagator; all particles are advanced on a global time step.
    \item \textbf{Self-gravitating MHD is untested.} Nothing prevents combining \texttt{--prop magneto-ve} with self-gravity, but neither the time-step control nor the Lorentz force have been validated in that regime.
    \item \textbf{Exact energy conservation of the cleaning terms assumes equal particle masses}, which all built-in test cases satisfy. But this might not be future proof. 
\end{itemize}

\newpage

\section{Details of code implementation}
Software documentation for the SPH-EXA team and future developers of the code. 

Need full velocity Jacobian $J_v$ for induction. Thus, overridden IAD driver. 

branching on curlv allocated? new neighbor lists. gradh correction.

hardcoded divB cleaning variables. $f_{clean}=\sigma_c = 1$. Would need courant limit if raising fclean. evolve psi/ch. Only conservative conjugate pair for equal particle masses. mj/mi resiudal factor. 

divBerr diagnostic unreliable on domains with B=0 regions/spots. 

Using SLR: \texttt{T du\_diss = diss\_op * dot(B\_ab\_raw, B\_ab);} instead of \texttt{norm2(B\_ab)}.
energy conjugate of dBdiss: dE\_B/dt goes as B\_ab . B'\_ab, not |B'\_ab|2. Identical for the non-reconstructing schemes (B'\_ab == B\_ab). The transverse part of the SLR correction does no work on the field, so charging heat for it would leak energy.

No block time step. self-gravitating MHD untested (timestep control, limit lorentz force movement)

Code \href{https://github.com/MuellerNico/sphexa/commits/develop/}{\textbf{GitHub fork}}
